% Generated draft: requires mathematical and source review.
\documentclass[11pt]{article}
\usepackage[margin=25mm]{geometry}
\usepackage{fontspec}
\setmainfont{Latin Modern Roman}
\usepackage{amsmath,amssymb,mathtools,unicode-math}
\setmathfont{Latin Modern Math}
\usepackage{enumitem,needspace,xurl,hyperref}
\setlength{\parindent}{0pt}
\setlength{\parskip}{5pt}
\setlength{\emergencystretch}{3em}
\newcommand{\N}{\mathbb{N}}
\newcommand{\Z}{\mathbb{Z}}
\newcommand{\Q}{\mathbb{Q}}
\newcommand{\R}{\mathbb{R}}
\newcommand{\C}{\mathbb{C}}
\newcommand{\F}{\mathbb{F}}
\newcommand{\A}{\mathbb{A}}
\newcommand{\PP}{\mathbb{P}}
\newcommand{\E}{\mathbb{E}}
\newcommand{\eps}{\varepsilon}
\newcommand{\dd}{\,\mathrm{d}}
\begin{document}
\section*{UnsolvedMath ID: 30006469}
\textbf{Source problem code:} catalogue entry
\subsection{Definitions and mathematical statement}
The source fixes an elliptic curve \(E\), but does not specify its ground field or any additional hypotheses on \(E\). Its notation presupposes reductions over prime fields. The following formulation is therefore conditional on having a specified good reduction \(E_p/\F_p\) at each rational prime under consideration; the intended class of elliptic curves remains unresolved.

Let \(\mathcal P_E\) denote the set of those rational primes \(p\) at which the specified reduction \(E_p/\F_p\) is good. For \(p\in\mathcal P_E\), define the Frobenius trace by
\[
a_p=p+1-\#E_p(\F_p),
\]
and let \(\overline{a_p}\) be its residue class modulo \(p\). Whenever \(p\nmid a_p\), define
\[
\operatorname{ord}_p(a_p)=\min\{n\in\Z_{\geq 1}:a_p^n\equiv 1\pmod p\}.
\]
The residue class \(\overline{a_p}\) is a primitive root modulo \(p\) if it generates \(\F_p^{\times}\), equivalently if \(\operatorname{ord}_p(a_p)=p-1\).

For the fixed curve \(E\), the two source questions are whether, respectively,
\[
\#\{p\in\mathcal P_E:p\nmid a_p,\ \operatorname{ord}_p(a_p)=p-1\}=\infty
\]
and
\[
\#\{p\in\mathcal P_E:p\nmid a_p,\ \operatorname{ord}_p(a_p)\geq 6\}=\infty.
\]
The source gives no restriction such as a specified ground field or a condition on the endomorphisms or torsion of \(E\). Whether further hypotheses were intended must be checked before these questions can be treated as a fully specified assertion about every curve in a particular class.

The original passage also asserts, as context, that if \(p\) does not divide the discriminant (or the conductor) of \(E\), and \(a_p\neq 0\), then the multiplicative order of \(a_p\) modulo \(p\) equals the degree of the extension of \(\F_p\) generated by the \(p\)-torsion points of \(E\) over \(\F_p\). This contextual assertion is supplied source material, not a verified auxiliary theorem here. Its precise formulation requires a ground field, a suitable discriminant or conductor, clarification of the meaning of the \(p\)-torsion points in characteristic \(p\), and an invertibility hypothesis: the integer condition \(a_p\neq 0\) alone does not express \(p\nmid a_p\). No small primes are silently discarded in the two displayed questions.
\subsection{Short English statement}
For a fixed elliptic curve with the prime-field reductions presupposed by the source, ask whether its Frobenius trace generates \(\F_p^{\times}\) for infinitely many good primes \(p\), and whether it has multiplicative order at least six for infinitely many such primes. The ground field and any additional curve hypotheses are unspecified. Supplied dated curation labels the pair partially solved and reports a sketch for the second question; this draft does not independently verify that assessment.
\subsection{Sources}
\begin{itemize}
\item Supplied local UnsolvedMath registry record, permanent ID \(30006469\), title: \emph{Primitive-Root Orders of Elliptic-Curve Frobenius Traces}; exact source problem code: \texttt{catalogue entry}. Collection: Oberwolfach Reports Open Problems. The record gives proposed year 2025 and names no proposer.
\item Supplied report citation: \emph{Analytic Number Theory (2026)}, Oberwolfach Reports, DOI 10.4171/owr/2025/51, pp.\ 2701--2764. Supplied source URL: \url{https://doi.org/10.4171/owr/2025/51}. The background describes Oberwolfach Reports as documenting mathematical workshops and their open-problem discussions. The original passage is described in the record as retained verbatim.
\item Supplied \texttt{literature\_sources} entry: ``Oberwolfach report 51/2025,'' \url{https://doi.org/10.4171/owr/2025/51}.
\item Literature assessment dated exactly \texttt{2026-08-22}: supplied dated curation, not independently verified. Its status is ``partially solved.'' It reports that the workshop report contains a sketched solution for the order-at-least-six question, but provides no published complete proof for that question or for the stronger primitive-root assertion, and consequently describes the pair as not fully settled. The background identifies this as dated literature triage derived from supplied OWR curation and calls for independent verification of time-sensitive status and cited claims.
\item The registry labels the statement \texttt{corrected\_verified}, with the qualification that only formula delimiters, notation, grammar, and typography were repaired and that mathematical content and scope were preserved. This supplied editorial label does not establish mathematical correctness. The citation's title year 2026 and the report designation 51/2025 are both retained as supplied; their relationship has not been checked.
\end{itemize}
\end{document}
